% Source: https://artofproblemsolving.com/community/c3h2530554_what_exactly_is_conditional_probability
% Stable identity: aops-topic-2530554-q01
\begin{enumerate}
\item[(a)] For events $A$ and $B$ with $P(B)>0$, the \emph{conditional probability} of $A$ given $B$ is defined by $P(A\mid B)=\dfrac{P(A\cap B)}{P(B)}$. Explain why, when all outcomes are equally likely, $P(A\mid B)$ is the fraction of the outcomes in $B$ that are also in $A$.
\item[(b)] Two fair dice are rolled. Let $A$ be the event that at least one die shows a $6$, and $B$ the event that the sum is $8$. Find $P(A)$, $P(A\mid B)$ and $P(B\mid A)$. Are $A$ and $B$ independent?
\end{enumerate}
